\clearpage\SourceLeaf{124}%
\pdfbookmark[1]{Der Zodiakus der Juno}{zodiacus1805}
\begin{center}
{\small DER ZODIAKUS DER JUNO.\par}
\medskip
{\small (Ein Nachtrag zum Aufsatz im Augustheft der \textit{M. C.} von 1804 [S.\,106 d.\,B.])\par}
\medskip\hrule\medskip
{\small Monatliche Correspondenz zur Beförderung der Erd- und Himmels-Kunde,\par
herausgegeben vom Freiherrn \textsc{von Zach}. März 1805.\par}
\medskip\hrule
\end{center}\medskip
Obwohl Juno bisher erst kurze Zeit beobachtet worden ist, scheinen
die Elemente ihrer Bahn schon genau genug zu sein, um die Grenzen
der Zone abzustecken, in der uns der Planet erscheinen kann. Das
dient dem Atlas, den man von seinem verdienstvollen Entdecker
erhoffen darf. Auf Bitten meines Freundes \textsc{Harding} habe ich
diese Arbeit umso lieber übernommen, als gerade die kleine,
lichtschwache Juno in weniger günstigen Lagen als in diesem Jahr von
allen drei neuen Planeten am schwersten zu beobachten sein wird und
daher detaillierte Sternkarten am nötigsten haben wird.

Nach meinen letzten Beobachtungen scheint mein vierter Satz von
Elementen für Juno noch so gut mit ihrer Bewegung übereinzustimmen,
dass ich die Bahn vorerst nicht zuverlässig weiter verbessern
könnte. Daher habe ich diese Elemente meinen Rechnungen zugrunde
gelegt und in der Schreibweise des erwähnten Aufsatzes die folgenden
Werte angesetzt.
\[
\begin{aligned}
e&=0.25684\\
k&=2.49630\\
g&=62^0\;19'\;35''\\
e'&=0.01679\\
k'&=0.999718\\
g'&=108^0\;30'\;1''
\end{aligned}
\]
\SourceLeaf{125}%
Daraus fand ich die folgende Bedingungsgleichung.
\[
\cos t'=0.400477\cos t-0.053103
\]

Die Werte von $\cos t'$ liegen also zwischen $+0.347374$ und
$-0.453580$. Die Werte von $t'$ liegen folglich einerseits zwischen
$69^0\;40'\;23''$ und $116^0\;58'\;25''$, andererseits zwischen
$234^0\;5'\;47''$ und $281^0\;23'\;49''$. Damit liegen die
heliozentrischen Positionen der Erde, von denen aus Juno auf ihren
Grenzen erscheinen kann, zwischen den Längen $240^0\;45'$ und
$288^0\;3'$ sowie zwischen $54^0\;6'$ und $101^0\;24'$. Die ersten
fallen etwa in die Zeit vom 22.\,Mai bis 10.\,Juli, in der Juno
einmal auf der nördlichen Grenze erscheinen muss. Die anderen fallen
in die Zeit vom 16.\,November bis 2.\,Januar, in der sie einmal die
südliche Grenze erreicht.

Auf dieser Grundlage wurden genau wie bei Ceres und Pallas $36$
Punkte auf der nördlichen und ebenso viele auf der südlichen Grenze
bestimmt. Zur unmittelbaren Berechnung der Rektaszensionen und
Deklinationen werden die folgenden Formeln und Konstanten verwendet.
\[
\begin{array}{rcrrr}
a&=&87^0&59'&22''\\
b&=&100&32&38\\
c&=&10&44&17\\
A&=&261&18&7\\
B&=&171&40&35\\
C&=&160&38&2
\end{array}
\]
\[
v=\text{wahre Anomalie der Juno}
\]
\[
\begin{aligned}
x&=\frac{\alpha\sin(v+323^0\;37'\;41'')}{1-e\cos v}\\[8pt]
y&=\frac{\ObservedBeta\sin(v+234^0\;0'\;9'')}{1-e\cos v}\\[8pt]
z&=\frac{\gamma\sin(v+222^0\;57'\;37'')}{1-e\cos v}
\end{aligned}
\]
\[
\begin{aligned}
\alpha&=k\sin a; &\log\alpha&=0.397030\\
\ObservedBeta&=k\sin b; &\log\ObservedBeta&=0.389902\\
\gamma&=k\sin c; &\log\gamma&=9.667556
\end{aligned}
\]

Zwischen diese $72$ Punkte wurden $144$ weitere eingefügt. Ihre
Rektaszensionen steigen von $5$ zu $5$ Grad. Die folgende Tabelle
enthält diese Punkte.
\SourceLeaf{126}
\begin{center}
\begingroup\small\setlength{\tabcolsep}{3pt}
\begin{tabular}{r|rrl|rrl||r|rrl|rrl}
\multicolumn{14}{c}{Zodiakus der Juno.}\\[4pt]
&\multicolumn{6}{c||}{Deklination}&&\multicolumn{6}{c}{Deklination}\\
AR.&\multicolumn{3}{c|}{\shortstack{der nördlichen\\Grenze}}&\multicolumn{3}{c||}{\shortstack{der südlichen\\Grenze}}&AR.&\multicolumn{3}{c|}{\shortstack{der nördlichen\\Grenze}}&\multicolumn{3}{c}{\shortstack{der südlichen\\Grenze}}\\\hline
$0^0$&$4^0$&$23'$&nördl.&$10^0$&$38'$&südl.&$180^0$&$6^0$&$59'$&nördl.&$3^0$&$9'$&südl.\\
$5$&$5$&$13$&&$10$&$8$&&$185$&$6$&$5$&&$3$&$47$&\\
$10$&$6$&$3$&&$9$&$36$&&$190$&$5$&$11$&&$4$&$27$&\\
$15$&$6$&$51$&&$9$&$3$&&$195$&$4$&$17$&&$5$&$8$&\\
$20$&$7$&$39$&&$8$&$29$&&$200$&$3$&$23$&&$5$&$50$&\\
$25$&$8$&$26$&&$7$&$54$&&$205$&$2$&$30$&&$6$&$32$&\\
$30$&$9$&$11$&&$7$&$19$&&$210$&$1$&$38$&&$7$&$15$&\\[3pt]
$35$&$9$&$54$&&$6$&$43$&&$215$&$0$&$49$&&$7$&$58$&\\
$40$&$10$&$34$&&$6$&$7$&&$220$&$0$&$11$&&$8$&$40$&\\
$45$&$11$&$13$&&$5$&$32$&&$225$&$0$&$44$&südl.&$9$&$21$&\\
$50$&$11$&$49$&&$4$&$56$&&$230$&$1$&$26$&&$10$&$0$&\\
$55$&$12$&$23$&&$4$&$21$&&$235$&$2$&$3$&&$10$&$38$&\\
$60$&$12$&$53$&&$3$&$47$&&$240$&$2$&$38$&&$11$&$13$&\\[3pt]
$65$&$13$&$21$&&$3$&$14$&&$245$&$3$&$8$&&$11$&$45$&\\
$70$&$13$&$45$&&$2$&$43$&&$250$&$3$&$34$&&$12$&$15$&\\
$75$&$14$&$7$&&$2$&$13$&&$255$&$3$&$55$&&$12$&$42$&\\
$80$&$14$&$25$&&$1$&$45$&&$260$&$4$&$12$&&$13$&$5$&\\
$85$&$14$&$40$&&$1$&$19$&&$265$&$4$&$24$&&$13$&$26$&\\
$90$&$14$&$51$&&$0$&$56$&&$270$&$4$&$32$&&$13$&$43$&\\[3pt]
$95$&$14$&$58$&&$0$&$35$&&$275$&$4$&$34$&&$13$&$56$&\\
$100$&$15$&$1$&&$0$&$17$&&$280$&$4$&$32$&&$14$&$7$&\\
$105$&$15$&$1$&&$0$&$2$&&$285$&$4$&$25$&&$14$&$14$&\\
$110$&$14$&$55$&&$0$&$10$&nördl.&$290$&$4$&$14$&&$14$&$17$&\\
$115$&$14$&$46$&&$0$&$19$&&$295$&$3$&$58$&&$14$&$18$&\\
$120$&$14$&$32$&&$0$&$24$&&$300$&$3$&$38$&&$14$&$16$&\\[3pt]
$125$&$14$&$14$&&$0$&$26$&&$305$&$3$&$14$&&$14$&$11$&\\
$130$&$13$&$52$&&$0$&$24$&&$310$&$2$&$45$&&$14$&$2$&\\
$135$&$13$&$26$&&$0$&$19$&&$315$&$2$&$14$&&$13$&$52$&\\
$140$&$12$&$56$&&$0$&$9$&&$320$&$1$&$38$&&$13$&$39$&\\
$145$&$12$&$22$&&$0$&$4$&südl.&$325$&$1$&$0$&&$13$&$24$&\\
$150$&$11$&$44$&&$0$&$20$&&$330$&$0$&$19$&&$13$&$6$&\\[3pt]
$155$&$11$&$3$&&$0$&$40$&&$335$&$0$&$25$&nördl.&$12$&$46$&\\
$160$&$10$&$19$&&$1$&$4$&&$340$&$1$&$9$&&$12$&$24$&\\
$165$&$9$&$32$&&$1$&$31$&&$345$&$1$&$57$&&$12$&$0$&\\
$170$&$8$&$43$&&$2$&$1$&&$350$&$2$&$45$&&$11$&$35$&\\
$175$&$7$&$52$&&$2$&$34$&&$355$&$3$&$34$&&$11$&$7$&\\
$180$&$6$&$59$&&$3$&$9$&&$360$&$4$&$23$&&$10$&$38$&\\
\end{tabular}
\endgroup
\end{center}
\begin{center}\rule{.25\linewidth}{.4pt}\end{center}
\clearpage
\thispagestyle{empty}
\noindent\SourceLeaf{127}\null\par
\clearpage
